\documentclass[tikz,border=3mm]{standalone}
\usepackage[T2A]{fontenc}
\usepackage[utf8]{inputenc}
\usepackage[russian]{babel}
\usepackage{amsmath}
\usepackage{xcolor}
\usetikzlibrary{arrows.meta,matrix}

\newcommand{\num}[1]{\mathtt{#1}}
\newcommand{\res}[1]{\textcolor{blue}{\num{#1}}}

\begin{document}
\begin{tikzpicture}[font=\rmfamily\Large]
\node[anchor=west, inner sep=0pt] (eq) {$
\begin{aligned}
\num{1006}:\num{8} &= \num{125}\ \text{остаток}\ \res{6} \\
\num{125}:\num{8} &= \num{15}\ \text{остаток}\ \res{5} \\
\num{15}:\num{8} &= \num{1}\ \text{остаток}\ \res{7} \\
\num{1}:\num{8} &= \num{0}\ \text{остаток}\ \res{1}
\end{aligned}
$};

\draw[-{Latex[length=3mm,width=2mm]}, red, line width=0.8pt]
  ([xshift=6mm]eq.south east) -- ([xshift=6mm]eq.north east);
\end{tikzpicture}
\end{document}
